\documentclass[a4paper, twocolumn, 11pt, twoside]{article}

% update your language here.
\usepackage[portuguese]{babel}
\usepackage{linguamatica-v14}

%% Please use UTF8 in your document



%% Set the bibliography style.
%% Styles provided for Spanish (Castilian), Catalan and Portuguese.
%% Other languages will be created on a required basis.
%% If you notice any error let us know so we can share a fixed one with others.

% \bibliographystyle{sp_por}   % This for Portuguese
% \bibliographystyle{sp_gal} % This for Galician
% \bibliographystyle{sp_esp} % This for Castilian
\bibliographystyle{sp_cat} % This for Catalan


%% You can leave this unchanged. 
%% It will be updated by the editors when your paper gets published.
\submitted{\OCT{} 2017}
\accepted{\DEC{} 2017}


%% Add your title in the main language used in the article
%% If it is too long, you can use  \\  to break whenever needed or adequate.
\title{Artigo de Exemplo da Linguamática}

%% The title in English is also mandatory
\titleEN{Linguamática Sample Article}


%% Add authors here.
%% First line with author name, email and ORCid when available (only ID, please!)
%% In cases where a second line is needed for the affiliation,
%% use the command \nl to separate them. 
\author{
  Alberto Simões~\email{asimoes@ipca.pt}~\orcid{0000-0001-6961-2660}
  \instituto{2Ai, School of Technology\nl Instituto Politécnico do Cávado e do Ave}   
  \and 
  Xavier Gómez Guinovart~\email{xgg@uvigo.es}
  \instituto{Universidade de Vigo}
  \and
  José João Almeida~\email{jj@di.uminho.pt}
  \instituto{Universidade do Minho}
}


\begin{document}
\maketitle

%% Add the abstract in the main language of the article.
%% If possible, do not add any citations here.
\begin{resumo}
  Este documento apresenta a estrutura e formatação básica de um artigo a ser submetido e publicado na Linguamática. O resumo poderá conter mais que um parágrafo, see assim for necessário. 

   Lorem ipsum dolor sit amet, consectetur adipiscing elit. Nunc eget feugiat augue. Nulla id molestie quam, accumsan lobortis arcu. Maecenas varius nunc at enim ornare, ac molestie nulla accumsan. Pellentesque ligula ipsum, maximus et venenatis vel, venenatis eu elit. Morbi eget rutrum nibh, nec rutrum augue. Pellentesque nec auctor metus. Mauris dapibus mi volutpat ipsum egestas, efficitur posuere odio ornare. Phasellus ut condimentum nunc. Etiam lobortis non mauris et aliquam. Donec vitae cursus sapien. Vestibulum consequat, felis sed tincidunt congue, purus ipsum sagittis magna, sit amet cursus nunc purus nec tortor. Sed odio felis, egestas id est at, varius venenatis nulla. Quisque auctor libero eget sapien rhoncus dapibus. Fusce et nisl eu purus elementum pretium vitae id purus. 
\end{resumo}

%% Add keywords in the article's main language, in lowercase
\palavraschave{exemplo de artigo, notas para autores}


%% Add the abstract in English
\begin{abstract}
  This document presents the structure and basic formatting for a Linguamática article.
   The abstract might contain more than one paragraph if needed.

   Lorem ipsum dolor sit amet, consectetur adipiscing elit. Nunc eget feugiat augue. Nulla id molestie quam, accumsan lobortis arcu. Maecenas varius nunc at enim ornare, ac molestie nulla accumsan. Pellentesque ligula ipsum, maximus et venenatis vel, venenatis eu elit. Morbi eget rutrum nibh, nec rutrum augue. Pellentesque nec auctor metus. Mauris dapibus mi volutpat ipsum egestas, efficitur posuere odio ornare. Phasellus ut condimentum nunc. Etiam lobortis non mauris et aliquam. Donec vitae cursus sapien. Vestibulum consequat, felis sed tincidunt congue, purus ipsum sagittis magna, sit amet cursus nunc purus nec tortor. Sed odio felis, egestas id est at, varius venenatis nulla. Quisque auctor libero eget sapien rhoncus dapibus. Fusce et nisl eu purus elementum pretium vitae id purus. 
\end{abstract}

%% Add the keywords in English
\keywords{article sample, notes for the authors}



\section{Introducción}

A formatação ao longo do documento é a normal em documentos \LaTeX, sem grandes alterações. No entanto,
algumas sugestões:
\begin{itemize}
\item Para dar \emph{ênfase} use sempre que possível o comando \verb.\emph.;
\item Para citar poderá usar o comando \verb.\citep. que cria referências entre
  parêntesis~\citep{Goodfellow-et-al-2016}. Para citar um \citet{Rumshisky-Stubbs-2017}, 
  use o comando \verb.\citet.;
\item Citações seguidas devem reaproveitar o comando de citação. Caso necessite de indicar a página
  a que se refere a citação, use \citep[p.~4]{Batista-2023}.
\item Ao criar entradas bibliográficas assegure-se da correção do seu
  conteúdo. Não abrevie nomes de autores. Não coloque os nomes dos
  editores de livros de atas. Não se esqueça dos números das páginas
  do documento.
\item Sempre que usar endereços web e outros tipos de URI, coloque-os com o comando \verb.\url. e,
  sempre que possível, em nota de fim de página.\footnote{Assim. \url{http://www.linguamatica.com}}
\item Nas notas de fim de página que sejam anexadas a palavras seguidas de pontuação, devem ser colocadas
  após a pontuação, como exemplificado no item anterior.
\item Tenha em atenção a diferença entre -, -- e ---. O primeiro será usado entre palavras, como em
  curto-circuito, o segundo em intervalos, como 10--20 ou PT--EN e o terceiro --- este --- para
  introduzir pequenos comentários.
\item As figuras devem ser legendadas e a legenda deve terminar com um sinal de pontuação.
\item As referências a figuras, tabelas ou secções devem ser criadas usando as ferramentas do \LaTeX{}.
\item Sempre que possivel garanta a qualidade das imagens importadas, usando PDF ou PNG.
\item Ao criar tabelas (Tabela~\ref{tab:1}) tente diminuir a quantidade de traços usada. Grande parte das
  tabelas são legíveis apenas com um par de linhas como demonstrado.
\end{itemize}

\begin{table}[htb]
  \centering
  \begin{tabular}{lrr}
    \toprule
    & \textbf{Homens} & \textbf{Mulheres} \\
    \midrule
    Crianças & 10~032 & 32~341 \\
    Adultos & 23~431 & 9~443 \\
    \bottomrule
  \end{tabular}
  \caption{Exemplo de tabela com poucos traços.}
  \label{tab:1}
\end{table}

\section{Other Title}

 Aliquam enim metus, laoreet vel luctus ut, accumsan eget felis. Nunc semper varius rutrum. Quisque posuere feugiat risus, sed cursus velit pulvinar eu. Etiam vitae cursus diam. Nulla cursus efficitur metus, ut tincidunt lorem finibus et \cite{Simoes-Guinovart-2019}. Vivamus pharetra ornare scelerisque. Mauris tempor magna non lectus ultricies, ut finibus dui bibendum. Curabitur quis auctor enim, vel pretium est. Donec a odio at justo commodo sollicitudin eu quis mi. Sed rhoncus nisi at bibendum volutpat. Integer est nunc, tincidunt et est ac, placerat porttitor odio \citep{Zhou-et-al-2024}. 

  Suspendisse magna tellus, interdum sit amet elit at, blandit lacinia est. Proin ullamcorper ligula fermentum urna lacinia, vitae faucibus urna porttitor. Cras sed nibh eleifend justo accumsan accumsan at quis ante. Donec vitae magna arcu. Vivamus luctus odio ut lorem pellentesque, in facilisis ante lobortis. Proin ut lacus tempus, finibus risus eget, vehicula est. Praesent ut blandit nunc. Nulla nec mi a dui scelerisque fringilla. Nam eros nunc, \cite{Almeida-2003} finibus sit amet ex sit amet, euismod rutrum est. Proin eget est at risus scelerisque lobortis. Aliquam placerat erat feugiat iaculis mattis. Curabitur suscipit nisi ligula, nec consectetur erat luctus in. Quisque quis ante purus. Sed tristique lacus ac convallis tincidunt. Quisque sit amet arcu imperdiet ipsum hendrerit luctus nec quis urna \citep{Simoes-2004}. 



\section*{Agradecimientos}

Os agradecimentos devem ser colocados sempre numa secção final, sem número, tal como
neste exemplo. Sempre que o autor assim o entender, poderá agradecer aos revisores.


\bibliography{references.bib}


\end{document}
